\documentclass[12pt,a4paper]{article}
\usepackage[utf8]{inputenc}
\usepackage[T1]{fontenc}
\usepackage{lmodern}
\usepackage[
    a4paper,
    left=1.7cm,
    right=1.7cm,
    top=1.8cm,
    bottom=1.7cm
]{geometry}
\usepackage{amsmath,amssymb}
\usepackage{microtype}
\usepackage{fancyhdr}
\usepackage[hidelinks]{hyperref}
\usepackage{titlesec}
\usepackage{enumitem}

\pagestyle{fancy}
\fancyhf{}
\lhead{\small Alexandros Sopasakis, 1968-04-08}
\rhead{\small KAW Mathematics Programme 2026}
\cfoot{\small\thepage}
\setlength{\headheight}{14pt}
\setlength{\parindent}{0pt}
\setlength{\parskip}{3pt}
\setlength{\abovedisplayskip}{5pt}
\setlength{\belowdisplayskip}{5pt}
\titleformat{\section}{\large\bfseries}{\thesection.}{0.5em}{}
\titleformat{\subsection}{\normalsize\bfseries}{\thesubsection}{0.5em}{}
\titlespacing*{\section}{0pt}{8pt}{3pt}
\titlespacing*{\subsection}{0pt}{6pt}{2pt}
\newcommand{\E}{\mathbb E}
\newcommand{\W}{W_1}
\newcommand{\piy}{\pi_y}
\newcommand{\pir}{\pi_{y,r}^{b}}
\newcommand{\epsb}{\varepsilon^{b}_{y}}
\newcommand{\epsy}{\varepsilon_{y}}
\newcommand{\normw}[1]{\left\|#1\right\|_w}
% Check the six-page submission route after flushing the final page.
\AtEndDocument{\clearpage\ifnum\value{page}>7
  \PackageError{ProposalLength}{Proposal exceeds six pages}{Shorten the text before submission.}
\fi}
\raggedbottom

\begin{document}
{\LARGE\bfseries When is local learning enough?\par}
\vspace{3pt}
{\large\bfseries Stability and sharp error bounds for conditional sampling on graphs\par}
\vspace{4pt}
Alexandros Sopasakis, Centre for Mathematical Sciences, Lund University

A field may contain a million coordinates while the research question concerns only one. \textit{How much} of the field must be read, and \textit{how accurately} must its local laws be learned? \\
I seek to recruit a postdoctoral mathematician to answer this for conditional sampling on graphs. Separately learned local scores (gradients of conditional log-densities) need not fit any joint law. A score which is accurate where the data lie can also drive a sampler into regions where it is wrong.

We ask: \emph{when does local learning suffice for accurate conditional sampling?} 
For a fixed target $S$, we seek $\W$-error of order $\varepsilon+e^{-\kappa r}$, up to polynomial factors in $r$, discretisation and finite-time errors, using only its radius-$r$ neighbourhood. Here  $\kappa$ is the spatial decay rate and $\varepsilon$ measures population score error, supplemented by explicit derivative control. 
Two aims drive the project: 
\begin{itemize}[nosep]
\setlength{\itemsep}{-10pt}
\setlength{\parskip}{0pt}
\setlength{\parsep}{0pt}
\setlength{\topsep}{0pt}
\setlength{\partopsep}{0pt}
\setlength{\leftmargin}{0em}
    \item[\textbf{(a)}] turn reference-law error into a local guarantee for nonconvex, incompatible learned updates, including transfer across an artificial boundary, and\\
\item[\textbf{(b)}]  choose radius, step size and run length from the desired accuracy. 
In a Gaussian benchmark, upper and lower truncation exponents agree when the full spatial resolvent is retained. 
\end{itemize}
The row-sum estimate needs an asymptotic radius about $19$ times larger at correlation $a=0.9$. Retaining this sharp exponent in the full learned bound is a further research objective.

\section{Research environment at Lund University}
\textbf{1.1 Host and mathematical setting.}
I am Senior Lecturer and Docent in Mathematics at Lund's Centre for Mathematical Sciences, where I have worked since 2011. With Katsoulakis and Plech\'a\v{c}, I established second-order weak error estimates for coarse-grained lattice dynamics, governed by coarse-graining scale relative to interaction range \cite{KPS06}. My work on stochastic closures \cite{KMS05}, interacting traffic models \cite{SK06} and lattice-free dynamics \cite{Sop12} further connects local interactions with macroscopic behaviour. With my doctoral student Donglin Liu, I developed GS-SINDy, combining Earth Mover's distance and group similarity with convergence analysis of its iteration \cite{LS25}. Our graph-enhanced diffusion work with Oskar \AA{}str\"om uses graph clustering and specialised decoders for spatiotemporal imputation \cite{LAS25}. 
The coarse-graining work asked this project's question for a fixed approximation: how a local error, measured against the interaction range, propagates through an interacting system. Here the approximation is learned.
%The proposed project connects these strands through local error analysis for learned stochastic updates.

I formulated the proposed project and carried out the preliminary work below: the Gaussian feedback example, the nonconvex numerics and the parameter example. With the fellow I will lead the boundary transfer (Theorem~A), closest to my coarse-graining error analysis, take part in the defect estimates and their certification, and connect the theory to our graph-imputation models \cite{LAS25}, which also supply the testbeds.
%I formulated this project and carried out the preliminary Gaussian calculations and nonconvex numerical exploration below. I will lead the boundary-transfer analysis with the fellow, contribute to the defect estimates and certification, and connect the theory to our graph-imputation models.

\textbf{1.2 Collaboration and supervision.}
We combine numerical analysis and mathematical statistics.\\
\noindent\textbf{(a)} \textit{Martin Friesen}, Senior Lecturer in Mathematical Statistics, has agreed to be second mentor. His work on ergodicity, concentration and interacting particle systems supports the nonconvex coupling analysis.

\noindent\textbf{(b)} \textit{Markos A.\ Katsoulakis}, Professor and Director of the Center for Applied Mathematics and Computation at UMass Amherst, has confirmed participation as co-mentor. We have eight joint papers, and his recent work links score-based models to Wasserstein proximal operators \cite{ZLLKO24}. He will contribute to analytical milestones, host a four-week visit in year 1 and visit Lund in year 2.

\noindent\textbf{(c)} \textit{Pablo Villanueva-Perez} (Synchrotron Radiation Research, Lund, ERC Starting Grant holder) has confirmed an optional MAX IV imaging study, connected to his sparse X-ray reconstruction work \cite{Zhang25}.
I have been main supervisor of Donglin Liu (PhD 2025) and Oskar \AA{}str\"om (2023--), co-supervisor of four further doctoral students, and supervisor or co-supervisor of three postdocs, including Shuzhen Li (main supervisor, 2025--). The fellow and I will meet weekly, with monthly meetings involving both co-mentors.

\section{Project and research area}
\textbf{2.1 What is known, and the missing result.} Dobrushin comparison propagates local discrepancies through spatial influence matrices \cite{Dob70,RvH14}. Sparse block samplers \cite{TMM20}, Wasserstein perturbation theory \cite{RS18} and reflection coupling \cite{Eberle16,MMS20} supply the analytical tools. Incompatible conditionals also arise in iterative imputation \cite{LGHSK14}. The open problem here is to control finite learned updates from reference-law errors while retaining spatial decay.
Cui, Liu and Tong \cite{CLT} bound marginal errors using target-law score discrepancies when the approximating density has the required locality. Their result compares scores of joint densities. Our learned drifts need not equal the score of their sampler's invariant law, so this route is not directly available. Gottwald and Reich \cite{GR26} establish stable localised samplers. Localised diffusion models \cite{GLMRT} analyse radius--statistical-error tradeoffs for reverse diffusion. Xi's comparison assumes weighted drift error along the learned trajectory \cite{Xi26}. We seek to derive that control from reference-law error and residual regularity. Langevin sampling can fail with arbitrarily small population $L^p$ score error \cite{CCSW26}, and numerical discretisation adds a stability issue \cite{Zho26}.
\textbf{The proposed novelty} is a spatial perturbation theorem, uniform in graph size, for the random-scan Euler sampler actually run, with nonconvex, incompatible drifts and an artificial boundary. Two estimates are new: a transfer bound in which the cut enters as an explicit forcing term propagated by the field's resolvent (Theorem~A), and a defect-feedback recursion that averages score error under the reference law, at the price of a controlled term $\Lambda d$, for drifts that are neither contractive nor compatible (Theorem~B). Verifiable block conditions give an abstract theorem, and a sinusoidal Gibbs family the first certified nonconvex example. %Dr Hyemin Gu is a potential candidate (Section~3).

%\textbf{The proposed advance} is a spatial perturbation theorem for the actual random-scan Euler sampler with nonconvex, incompatible drifts and an artificial boundary. The new estimates will transfer reference-law errors into bounds on local observables, uniformly in graph size. Verifiable block conditions will give an abstract theorem, with a sinusoidal Gibbs family as the first certified nonconvex example. Dr Hyemin Gu is a potential candidate (Section~3).

\textbf{2.2 Sampler, assumptions and a baseline.}
Let $X$ have a smooth, positive, finite-range Gibbs density on a bounded-degree graph with $N$ coordinates. Given $X_O=y$, write $\piy=\mathcal L(X_U\mid X_O=y)$ and partition $U$ into bounded blocks $B_i$. The target $S$ is a fixed union of blocks, identified with their indices in sums. Retain its radius-$r$ neighbourhood $I=I_r(S)$ in the block dependency graph and clamp the remaining unobserved coordinates to $b$, obtaining $\pir$. Observations remain fixed. A learned drift $s_i(u,z,y)$ reads the local boundary $z$ and observations available in the retained neighbourhood. Each outer update chooses one of the $k_r=|I|$ blocks uniformly and takes $n$ steps
\begin{equation}\label{eq:update}
 Z_{\ell+1}=Z_\ell+h s_i(Z_\ell,z,y)+\sqrt{2h}\,\xi_{\ell+1},
 \qquad \xi_\ell\sim N(0,I_{|B_i|}),\qquad \tau=nh.
\end{equation}
Write $\mu_T$ for the resulting law after $T$ updates. On $S$, $\W$ uses the sum of block Euclidean distances. The inputs are population losses
\[
 \varepsilon_{y,j}^2=\E_{\piy}\|s_j-s_j^\star\|^2,
 \qquad \varepsilon^b_{y,j}=\E_{\pir}\|s_j-s_j^\star\|,
\]
where $s_j^\star$ is the true conditional score, evaluated with the appropriate boundary. Finite-sample generalisation is outside this project. These losses concern the specified $y$. An unconditional training loss does not ensure accuracy for rare observations.

For a nonnegative matrix $A$, set $\normw{A}=\sup_i w_i^{-1}\sum_j A_{ij}w_j$, with $0<w_-\leq w_i\leq w_+$ uniformly in $N$. Let $C$ bound the $\W$ influences of the full conditional specification of $\piy$, with $\normw{C}\leq q<1$ uniformly over the bounded observations considered. Assume uniformly bounded block dimensions, finite interaction range, uniform local Lipschitz and linear-growth constants for the true scores, and uniform fourth moments per block under the full and clamped reference laws.

\textbf{Baseline proposition (months 1--4).} If learned block drifts are $L$-Lipschitz, $\widetilde m$-contractive in their own variable and have cross-block Lipschitz matrix $a$, then a synchronous comparison gives the intended bound
\begin{equation}\label{eq:baseline}
 \W((\mu_T)_S,(\pir)_S)
 \leq\sum_{i\in S}\left[(I-A_h)^{-1}
 \frac{h(\epsb+c\sqrt h\,\mathbf1)}{1-\rho_h}\right]_i
 +K_SM_0e^{-\gamma_h T/k_r}.
\end{equation}
Here $\rho_h=(1-2\widetilde mh+L^2h^2)^{1/2}<1$, $A_h=ha/(1-\rho_h)$ and $\gamma_h=(1-\rho_h^n)(1-\normw{A_h})$. We require $\normw{A_h}\leq1-\zeta$ for fixed $\zeta>0$. The local initial discrepancy is $M_0=\max_{i\in I}w_i^{-1}\E\|\widehat X_{0,i}-X_{0,i}\|$, with $X_0\sim\pir$. This baseline yields a unique invariant law and tests the error decomposition before removing learned-drift contractivity.

\textbf{2.3 First main result: transfer across an artificial boundary.}
Clamping changes the distribution on which the learned residual is tested. We will quantify that change before inserting score errors into a sampling bound. Set $D_I=(I-C_{II})^{-1}$ and $m_j^b=\E_{\piy}\|X_{B_j}-b_j\|$ for $j\notin I$. A spatial comparison coupling is to control the interior discrepancy by
\[
 t_I=D_I C_{I,I^c}m_{I^c}^b,\qquad t_{I^c}=m_{I^c}^b.
\]
If $L_{jk}$ are blockwise Lipschitz constants of $\|s_j-s_j^\star\|$, the coupling suggests $\varepsilon^b_{y,j}\leq\E_{\piy}\|s_j-s_j^\star\|+\sum_kL_{jk}t_k\leq\varepsilon_{y,j}+\sum_kL_{jk}t_k$, the last step by Cauchy--Schwarz.
%If $L_{jk}$ are blockwise Lipschitz constants of $\|s_j-s_j^\star\|$, the coupling suggests $\varepsilon^b_{y,j}\leq\varepsilon_{y,j}+\sum_kL_{jk}t_k$. The cut enters as an explicit forcing term involving the discarded coordinates.

\textbf{Target Theorem A.} For the nonnegative sampling resolvent $R$ arising in \eqref{eq:baseline} (or Thm~B), establish
\begin{equation}\label{eq:transfer}
 \sum_{i\in S}[R\epsb]_i
 \leq\sum_{i\in S}[R\epsy]_i+\sum_{i\in S}[RLt]_i,
 \qquad \W((\pir)_S,(\piy)_S)\leq\sum_{i\in S}t_i.
\end{equation}
The substantive target is an explicit decaying bound on the \emph{weighted} correction $RLt$, uniform in $N$ and $r$. We will retain the two resolvents to distinguish the correlation range of the field from the propagation range of sampler errors. Under polynomial neighbourhood growth and suitable exponential kernel bounds, we aim for $K_SM_b(1+r)^p e^{-\kappa r}$, where $M_b=\sup_{j\notin I}m_j^b$ and $p,\kappa$ are explicit. On general bounded-degree graphs we will state the actual weighted tail bound, rather than assume polynomial volume growth. This is the bridge from a population loss to a local algorithmic guarantee.

\textbf{2.4 Main challenge: nonconvex and incompatible learned updates.}
A drift may agree with $-u$ on $[-R,R]$ and point outward beyond $2R$, while its Gaussian-weighted error tends to zero. Additional structure is therefore necessary. Our first test family has
\begin{equation}\label{eq:potential}
 V(x)=\sum_i\left[\frac{\alpha x_i^2}{2}+\eta(1-\cos x_i)\right]
 -\sum_{\{i,j\}\in E}J_{ij}x_ix_j,
 \quad \eta>\alpha>0,\quad\sup_i\sum_j|J_{ij}|<\alpha,
\end{equation}
with symmetric finite-range $J$. Its score is globally Lipschitz and its tails are quadratically confined, but every one-site potential is nonconvex near $\pm\pi$. A separate, stronger mixing condition will be imposed. The learned drift is
\[
 s_i(u,z)=-\alpha u+\sum_jJ_{ij}z_j+g_i(u,z),\qquad
 g_i^\star(u)=-\eta\sin u.
\]
For fixed $B\geq\eta$, let $\mathcal G$ consist of local $C^1$ functions satisfying $|g_i|\leq B$, $\|\partial_u(g_i-g_i^\star)\|_\infty\leq\lambda_1$ and $\|\partial_{z_j}g_i\|_\infty\leq\lambda_2$ for each neighbour $j$. Degree factors are explicit. If $\lambda_1<\eta-\alpha$, every drift in this class fails global one-site contractivity. For $B>\eta$, residuals $g_i-g_i^\star=a_i+\sum_j b_{ij}\tanh z_j$ with $|a_i|+\sum_j|b_{ij}|\leq B-\eta$ and $|b_{ij}|\leq\lambda_2$ 
give explicit incompatible examples with $\lambda_1=0$ when $b_{ij}$ is nonsymmetric, since the residual does not depend on $u$.
%give explicit incompatible examples when $b_{ij}$ is nonsymmetric. They have $\lambda_1=0$.

The derivative bound is a global accuracy assumption additional to score loss. Bounded architectures can enforce amplitude and cross-derivative bounds. Certifying derivative accuracy relative to an unknown score is harder. We will test certificates on controlled models and distinguish assumptions from quantities observable in training. The explicit potential, with known score, is a proof and validation model. We will pursue broader potentials through the block criterion, without assuming arbitrary multimodality tractable.

\textbf{Local defect mechanism.} Let $P_i$ be the exact one-block diffusion over time $\tau$, which preserves $\pir$ when the starting joint law is $\pir$. In a common local metric $\mathfrak d$ with $c_-|u-v|\leq\mathfrak d(u,v)\leq c_+|u-v|$, seek initial-state contraction $\theta_i<1$ and boundary influences $M_{ij}$. Glue a coupling of learned and exact updates from one state to a coupling of exact updates from different states. Reflection coupling supplies the latter for nonconvex blocks. If the one-update defect satisfies $b_i\leq\widetilde b_i+e_i$, with block Lipschitz constants $\Lambda_{ij}$ for $\widetilde b_i$, then coupling to the stationary reference gives
\[
 \E b_i(\widehat X)\leq\bar b_i+\sum_j\Lambda_{ij}d_j+c_2\sqrt h,
 \qquad \bar b_i=\E_{\pir}\widetilde b_i,
\]
where $d_j=\E\mathfrak d(\widehat X_j,X_j)$. Thus the defect can be averaged under the reference law at the price of a controlled feedback term $\Lambda d$.

For the scalar family, synchronous comparison gives a candidate score-defect envelope
\begin{equation}\label{eq:defect}
 \widetilde b_i(x)=\int_0^\tau e^{\ell(\tau-t)}\E_x|r_i(Y_t)|\,dt,
 \quad r_i=g_i-g_i^\star,\quad \ell=\eta-\alpha+\lambda_1,
\end{equation}
in Euclidean distance, where $Y$ is the exact block diffusion with frozen boundary. Stationarity gives $\bar b_i\leq (e^{\ell\tau}-1)\varepsilon^b_{y,i}/\ell$. Flow sensitivity should yield $\Lambda_{ii}\leq C(\tau)\lambda_1$ and $\Lambda_{ij}\leq C(\tau)(\lambda_2+\lambda_1|J_{ij}|)$. Metric-equivalence factors must be included when using reflection distances. A constant residual has a constant envelope even when its exact update defect varies with the starting point. The confining linear part and bounded $g_i$ will provide fourth-moment bounds, uniform over outer time and graph size, to control the Euler defect $\E e_i(\widehat X_t)\leq c_2\sqrt h$ throughout the transient.

\textbf{Target Theorem B.} Fix $\alpha,\eta,B$, the degree bound and an admissible $\tau$. Prove explicit positive thresholds for $h$, $\sup_i\sum_j|J_{ij}|$, $\lambda_1$ and $\lambda_2$ for which, uniformly over $\mathcal G$,
\[
 H=I-\operatorname{diag}(\theta_i),\qquad
 F=H^{-1}(M+\Lambda),\qquad \normw{F}\leq1-\zeta,
 \qquad H^{-1}\bar b\leq c_1\epsb.
\]
The random-scan recursion should then imply
\begin{equation}\label{eq:main}
 \W((\mu_T)_S,(\pir)_S)
 \leq c_-^{-1}\sum_{i\in S}\left[(I-F)^{-1}
 (c_1\epsb+c_2'\sqrt h\,\mathbf1)\right]_i
 +K_SM_0 e^{-\gamma T/k_r},
\end{equation}
with $M_0$ defined in $\mathfrak d$ and $\gamma=\min_i\{1-\theta_i-w_i^{-1}\sum_j(M_{ij}+\Lambda_{ij})w_j\}>0$. Constants depend on the structural bounds and target size, not on $N$ or $r$. This controls target error from a bounded-moment start. Mixing of the learned chain requires a separate argument. Combining Theorems~A and~B is the central deliverable.

\textbf{Preliminary evidence and what it does not prove.} For two Gaussian sites with scores $-x_i+cx_j$, perturb the cross coefficients to $c+\delta$ and $c-\delta$, where $0\leq c<1$ and $0<|\delta|<1-c$. Exact inner integration gives $\theta=e^{-\tau}$, off-diagonal entries $c$ in $H^{-1}M$ and $|\delta|$ in $H^{-1}\Lambda$. The stationary coupled discrepancy obeys $d_i\leq\sqrt{2/\pi}\,\varepsilon/(1-c-|\delta|)$, where $\varepsilon=|\delta|/\sqrt{1-c^2}$. This verifies the feedback mechanism with incompatible drifts for every $\tau>0$.

For \eqref{eq:potential} with $\alpha=1$, $\eta=1.5$, $\tau=4$, my generator discretisation on $[-14,14]$ (561 points) estimated $\theta\approx0.39$ and sensitivity to the scalar boundary field $\approx2.01$, searching $|x|\leq10$ and 16 fields. Refinement to 801 points and 24 fields was stable, and the code recovered the Gaussian constants at $\eta=0$. Quadrature gives the separate equilibrium criterion $\sum_j|J_{ij}|<1/\sup_f\operatorname{Var}_{\pi_f}U\approx0.393$. These are numerical estimates, not certified global bounds.
On a nearest-neighbour chain with $J_{ij}=0.05$ and $\lambda_1=\lambda_2=0.005$, the preliminary constants give $\normw{F}\approx0.81$ and $c_1\approx21$. The resulting score multiplier is about $110$, so admissibility is not yet a practically tight guarantee. We will optimise $\tau$ and the defect envelope while certifying the margin, combining analytic tails, validated discretisation and bounds between grid points. Field reduction must respect the joint shift $(u,f)\mapsto(u+2\pi,f+2\pi\alpha)$. For $\tau\leq1$, the Euclidean estimate exceeded one. We will then use reflection distances, accounting for their equivalence constants.

\textbf{2.5 Sharp spatial benchmark and computational consequence.}
For the stationary chain $X_{j+1}=aX_j+\sqrt{1-a^2}\,\xi_{j+1}$, $0<a<1$, observations $X_d=\pm1$ give laws $\pi_\pm=N(\pm a^d,1-a^{2d})$ at the origin. A sampler unable to read $X_d$ has the same output law $\nu$ in both cases, hence
\[
 \max\{\W(\nu,\pi_+),\W(\nu,\pi_-)\}\geq a^d.
\]
For Gaussian precision $Q$ with nonpositive off-diagonal entries, $C=I-D_Q^{-1}Q$ and $(I-C)^{-1}=\Sigma D_Q$, where $D_Q=\operatorname{diag}(Q)$ \cite{MJW06}. On the infinite chain these weights are $(1+a^2)a^{|i-j|}/(1-a^2)$, matching the lower spatial exponent. The row-sum estimate $q^r$, $q=2a/(1+a^2)$, requires an asymptotically larger radius by $\log(1/a)/\log(1/q)$: about $19$ at $a=0.9$ and $2\xi$ near criticality, with $\xi=1/\log(1/a)$. This establishes sharpness for truncation, not for the entire learned error or runtime. Preserving the exponent after residual feedback and boundary transfer is an ambitious further target.
For fixed structural margins, bounded $M_0$, fixed $\tau=nh$ and local evaluation cost, an error floor of order $\varepsilon+\sqrt h$ gives a score-evaluation budget of order $|I_r(S)|\delta^{-2}\log(1/\delta)$ for accuracy $\delta$, with $h\asymp\delta^2$, sufficient score accuracy, and radius chosen from the spatial bound. This online cost excludes training and certification. Dependence on correlation length, graph growth and stability margins will be explicit.

\section{Candidate and planned fellowship}
\textbf{3.1 Recruitment, independence and contribution to Lund.}
Dr Hyemin Gu has expressed interest in applying. She defended her PhD at UMass Amherst in September 2025 with Katsoulakis and is now a postdoc there. Her first-author work develops Lipschitz-regularised gradient flows and generative particle algorithms for scarce high-dimensional data \cite{GBPRK24}. Her recent first-author preprint studies sharp stability thresholds and certification for residual architectures \cite{GTSK26}. Transport-based analysis and stability of learned dynamics are the two skills this project needs. At Lund she would develop a distinct line: spatial perturbation theory for interacting learned samplers. Friesen adds coupling and ergodicity expertise, and my group supplies stochastic error analysis and graph-based modelling. Katsoulakis provides continuity and specialist input. The fellow will lead Theorem~B, choose its extension and shape the research programme at Lund. Authorship will reflect contributions, with freedom to publish independently. The position will be advertised internationally and filled in open competition, with a committee comprising myself, Friesen and Villanueva-Perez. We seek a recent PhD in probability, stochastic analysis, optimal transport or numerical analysis, with evidence of independent mathematical work.

\textbf{3.2 Two-year programme.} Starting 1 July 2027. \textbf{Months 1--4: foundations and early stress tests.} Establish the baseline and cost accounting, verifying every term on Gaussian examples. Reproduce the nonconvex estimates and initiate certification immediately. By month~4, identify which constants need improvement and fix the metrics and moment assumptions.
\textbf{Months 3--10: spatial transfer and finite updates.} Prove Theorem~A and the finite-time recursion, and close the convex-block case. The first manuscript will combine weighted boundary transfer with incompatible finite updates and an explicit computational guarantee. During the Amherst visit, test the defect analysis against alternative perturbation arguments. By month~8, decide whether longer-time Euclidean contraction or reflection distances provide the stronger certified regime.
\textbf{Months 8--20: the nonconvex theorem.} Prove Theorem~B for the explicit family, including transient moments, metric constants and global numerical certification. Pursue broader one-site potentials through the block criterion. The second manuscript will give a certified nonconvex regime and quantify deterioration of the error bound near its boundary. At month~12 the fellow chooses an extension: signed analysis for frustrated Gaussian models, or guarantees for typical observations using disagreement-percolation ideas \cite{vdBM94}. The aim is an independent result beyond the initial theorem.
\textbf{Months 18--24: testbeds and a continuing research line.} Compare certified bounds with observed local errors on a non-Gaussian field and a sensor network, and release reproducible code. The optional phantom/MAX IV study will test the consequences of violated locality. The fellow will run a working seminar on sampling and localisation, organise a two-day national workshop in year~2 and prepare a Swedish Research Council starting-grant application, subject to eligibility.

\textbf{3.3. Risks and lasting value.}
The principal risks are a weak contraction margin, excessive defect amplification and slow decay after transfer. Each milestone will report both admissibility and the full accuracy multiplier. If the sinusoidal example needs better constants, we will improve the envelopes or change metric while preserving the nonconvex goal. If transfer fails, we will identify the missing assumption through a counterexample and prove the corresponding conditional result. The convex fallback must still deliver spatially weighted population-to-sampler bounds for incompatible finite updates and artificial boundaries. Reproducing a standard perturbation bound alone will not count as success.
The lasting contribution to Lund is a mathematical research line connecting probability with learned stochastic dynamics, developed by a fellow with ownership of its main theorem and extensions. We will invite Jimmy Olsson (KTH), Fredrik Lindsten (Link\"oping) and Moritz Schauer (Chalmers) to the seminar and workshop, and seek participation in Mittag-Leffler's \emph{Frontiers in Optimal Control: Geometry, Complexity, and Learning} programme in autumn~2028. 
%The ambition is both a rigorous answer to when locality is enough and an independent researcher equipped to push that answer beyond the first admissible class.
%\vspace{-.2cm}
\begingroup
\footnotesize
\begin{thebibliography}{99}
\setlength{\itemsep}{0pt}\setlength{\parskip}{0pt}
\bibitem{KPS06} M.~A. Katsoulakis, P. Plech\'a\v{c} and A. Sopasakis. \href{https://doi.org/10.1137/050637339}{Error analysis of coarse-graining for stochastic lattice dynamics}. \emph{SIAM J. Numer. Anal.} 44, 2270--2296 (2006).
\bibitem{KMS05} M.~A. Katsoulakis, A.~J. Majda and A. Sopasakis. \href{https://doi.org/10.4310/CMS.2005.v3.n3.a9}{Multiscale couplings in prototype hybrid deterministic/stochastic systems. Part II: Stochastic closures}. \emph{Commun. Math. Sci.} 3, 453--478 (2005).
\bibitem{SK06} A. Sopasakis and M.~A. Katsoulakis. \href{https://doi.org/10.1137/040617790}{Stochastic modeling and simulation of traffic flow: asymmetric single exclusion process with Arrhenius look-ahead dynamics}. \emph{SIAM J. Appl. Math.} 66, 921--944 (2006).
\bibitem{Sop12} A. Sopasakis. \href{https://doi.org/10.4208/cicp.110211.200611a}{Lattice free stochastic dynamics}. \emph{Commun. Comput. Phys.} 12, 691--702 (2012).
\bibitem{LS25} D. Liu and A. Sopasakis. \href{https://doi.org/10.1063/5.0214404}{Enhancing sparse identification of nonlinear dynamics with Earth-Mover distance and group similarity}. \emph{Chaos} 35, 033139 (2025).
\bibitem{LAS25} D. Liu, O. \AA{}str\"om and A. Sopasakis. Graph-enhanced diffusion models for spatiotemporal imputation. Under revision, \emph{Integrated Computer-Aided Engineering}. Paper IV, D. Liu, \emph{System Identification and Data-Driven Modeling}, Lund PhD thesis (2025).
\bibitem{ZLLKO24} B.~J. Zhang, S. Liu, W. Li, M.~A. Katsoulakis and S. Osher. \href{https://doi.org/10.1137/24M1644584}{Wasserstein proximal operators describe score-based generative models and resolve memorization}. \emph{SIAM J. Math. Data Sci.} 8, 623--648 (2026).
\bibitem{Zhang25} Y. Zhang, Z. Yao, R. Kl\"ofkorn, T. Ritschel and P. Villanueva-Perez. \href{https://doi.org/10.1038/s44172-025-00390-w}{4D-ONIX for reconstructing 3D movies from sparse X-ray projections via deep learning}. \emph{Commun. Eng.} 4, 54 (2025).
\bibitem{Dob70} R.~L. Dobrushin. \href{https://doi.org/10.1137/1115049}{Prescribing a system of random variables by conditional distributions}. \emph{Theory Probab. Appl.} 15, 458--486 (1970).
\bibitem{RvH14} P. Rebeschini and R. van Handel. \href{https://doi.org/10.1007/s10955-014-1087-7}{Comparison theorems for Gibbs measures}. \emph{J. Stat. Phys.} 157, 234--281 (2014).
\bibitem{TMM20} X.~T. Tong, M. Morzfeld and Y.~M. Marzouk. \href{https://doi.org/10.1137/19M1284014}{MALA-within-Gibbs samplers for high-dimensional distributions with sparse conditional structure}. \emph{SIAM J. Sci. Comput.} 42, A1765--A1788 (2020).
\bibitem{RS18} D. Rudolf and N. Schweizer. \href{https://doi.org/10.3150/17-BEJ938}{Perturbation theory for Markov chains via Wasserstein distance}. \emph{Bernoulli} 24, 2610--2639 (2018).
\bibitem{Eberle16} A. Eberle. \href{https://doi.org/10.1007/s00440-015-0673-1}{Reflection couplings and contraction rates for diffusions}. \emph{Probab. Theory Relat. Fields} 166, 851--886 (2016).
\bibitem{MMS20} M.~B. Majka, A. Mijatovi\'c and \L. Szpruch. \href{https://doi.org/10.1214/19-AAP1535}{Nonasymptotic bounds for sampling algorithms without log-concavity}. \emph{Ann. Appl. Probab.} 30, 1534--1581 (2020).
\bibitem{LGHSK14} J. Liu, A. Gelman, J. Hill, Y.-S. Su and J. Kropko. \href{https://doi.org/10.1093/biomet/ast044}{On the stationary distribution of iterative imputations}. \emph{Biometrika} 101, 155--173 (2014).
\bibitem{CLT} T. Cui, S. Liu and X.~T. Tong. Stein's method for marginals on large graphical models. \href{https://arxiv.org/abs/2410.11771}{arXiv:2410.11771}, revised 2026.
\bibitem{GR26} G.~A. Gottwald and S. Reich. \href{https://doi.org/10.1016/j.jcp.2025.114583}{Localized Schr\"odinger bridge sampler}. \emph{J. Comput. Phys.} 548, 114583 (2026).
\bibitem{GLMRT} G.~A. Gottwald, S. Liu, Y. Marzouk, S. Reich and X.~T. Tong. Localized diffusion models. \href{https://arxiv.org/abs/2505.04417}{arXiv:2505.04417}, revised 2026.
\bibitem{Xi26} W. Xi. When do local score models extrapolate across size? A diagnostic theory and benchmark. \href{https://arxiv.org/abs/2606.09705}{arXiv:2606.09705} (2026).
\bibitem{CCSW26} D.~Y. Cao, A.~Y. Chen, K. Sridharan and Y. Wu. On the robustness of Langevin dynamics to score function error. ICML 2026. \href{https://arxiv.org/abs/2603.11319}{arXiv:2603.11319}.
\bibitem{Zho26} Y. Zhou. Score accuracy along the forward diffusion does not certify numerical stability in diffusion sampling. \href{https://arxiv.org/abs/2607.08757}{arXiv:2607.08757} (2026).
\bibitem{MJW06} D.~M. Malioutov, J.~K. Johnson and A.~S. Willsky. Walk-sums and belief propagation in Gaussian graphical models. \emph{J. Mach. Learn. Res.} 7, 2031--2064 (2006).
\bibitem{GBPRK24} H. Gu, P. Birmpa, Y. Pantazis, L. Rey-Bellet and M.~A. Katsoulakis. \href{https://doi.org/10.1137/23M1587841}{Lipschitz-regularized gradient flows and generative particle algorithms for high-dimensional scarce data}. \emph{SIAM J. Math. Data Sci.} 6, 1205--1235 (2024).
\bibitem{GTSK26} H. Gu, M. Tyrrell, T. Sahai and M.~A. Katsoulakis. Certifying residual architectures from their primitives: a sharp stability threshold. \href{https://arxiv.org/abs/2607.14576}{arXiv:2607.14576v2} (2026).
\bibitem{vdBM94} J. van den Berg and C. Maes. \href{https://doi.org/10.1214/aop/1176988728}{Disagreement percolation in the study of Markov fields}. \emph{Ann. Probab.} 22, 749--763 (1994).
\end{thebibliography}
\endgroup
\end{document}
